% Pursuing Stacks (cotes 134-2 to 134-8): archive cote 134-6, site page 15 (PDF page 16), typescript page 334.
% typescript. Draft compared: ps5.tex:2103-2184.
% AI-transcribed by Claude Opus 5.5 (claude-opus-5-5, adaptive thinking, effort high), Message Batches API, batch ps3 (msgbatch_01XdFmuVVBAcgc7EFxdtCqsj), 2026-11-12; not yet checked by eye.
\scanpage{134-6/15}
\mspage{334}

This being \xout{\illegible{1}} clear,  the non-zero coefficients which occur in this
linear combination are all $\pm$ 1 ,  and moreover the maps which target a
\[ b \hand{\longrightarrow} a \qquad (b \quad A_{n-1}) \]
which occur with non-zero coefficient are exactly \uline{all monomorphisms}
from objects b in $A_{n-1}$ into a.  Thus,  the differential operators are
known,  when we know ,  for a\over{\unsure{ny}}{ll} monomorphisms in A
(13) \qquad $\hand{\partial}: b \hand{\to} a$ \quad , \quad with \quad dim a = dim b + 1
(the so-called ``\uline{boundary maps}''),  the corresponding coefficient\hand{s}
(14) \qquad $\hand{\varepsilon_{\partial} \in \{\pm 1\}}$ .
Instinct tells us,  at this point,  that we may get into trouble,  when
trying to define (in a more or less general case)  boundary operations
in such a way,  because of the ambiguity in the definition of subobjects\hand{,}
namely, \deleted{by the} \inserted{because of possible} \uline{existence of isomorphisms which may not be identities.}
But precisely,  in the standard cases we are \xout{\illegible{1}} copying from,  \uline{any}
\uline{isomorphism is an identity}\hand{!}

d)  For describing the ``signatures'' (14),  in one of the ``standard''
cases,  we still need to remark that  for any a in A,  we have
(15) \qquad i(a)  is an n-cell\xout{x},  with  n = dim (a) \quad ,
and the choice of the \xout{signs} signatures will be determined by a choice
of orientations
(16) \qquad $\hand{\omega}_a$ an \uline{orientation} of the n-cell $\hand{|}i(a)\hand{|}$ \quad ,
(a notion which could be given a purely combinatorial definition,  by
induction on the dimension of a given ordered set whose geometrical
realization is a va\unsure{t|r}iety ...).  We the\hand{n} get the\deleted{n} ``\uline{Stokes rule}''
(17) \quad For a boundary map $\hand{\partial}: b \hand{\to} a$ \quad , $\hand{\varepsilon_{\partial}} = +1$  iff $\hand{\omega}_b$ is ``induced'' à
la Stokes by $\hand{\omega}_a$\hand{,} via the induced orientation on the boundary
of  i(a) \hand{(} which is the union of the images of all  i(b) 's ,
for all boundary operations with target a\hand{)}.

Whether or not we are in a ``standard'' case,  if A is any category,
such that for any object a of A,  the ordered set i(a)  of it's subob-%
jects in A  is finite,  and it's geometrical realization is an n-cell
( call n the ``dimension'' of a) ,  and if moreover for a given n,  the
set $A_n$ of objects with dimension n is finite,  and also (to be on the
safe side !)  \deleted{we} assum\over{e}{\unsure{ing}} that all isomorphisms are identities,  then
for \uline{any} choice  of orientations (16),  giving rise to a system of

\margin{15}
